\documentclass[12pt]{article}
\usepackage[T2A]{fontenc}   
\usepackage[utf8]{inputenc}
\usepackage[russian]{babel}
\usepackage{amsthm, amsmath, amssymb}
\usepackage{graphicx}
\usepackage{float}
\usepackage{multirow}
\usepackage{enumitem}
\usepackage{listings}
\usepackage{xcolor}
\usepackage{placeins}
\usepackage{titlesec}
\usepackage{caption}
\usepackage{tikz}
\usepackage{setspace}
\usepackage[a4paper, mag=1000,
    left=25mm, right=25mm, top=18mm, bottom=20mm,
    bindingoffset=0mm,
    footskip=3.5ex,
    headheight=0cm, headsep=0cm
]{geometry}
\usetikzlibrary{arrows.meta, positioning, calc}
\usepackage{pgfplots}
\pgfplotsset{compat=1.18} % Настройки размера и совместимости

% Настройки для красивого отображения кода Python
\definecolor{codegreen}{rgb}{0,0.6,0}
\definecolor{codegray}{rgb}{0.5,0.5,0.5}
\definecolor{codepurple}{rgb}{0.58,0,0.82}
\definecolor{backcolour}{rgb}{0.95,0.95,0.92}


\lstset{
    language=Python,
    basicstyle=\ttfamily\small,
    keywordstyle=\color{blue},
    commentstyle=\color{green!50!black},
    stringstyle=\color{red},
    showstringspaces=false,
    breaklines=true,
    frame=single,
    numbers=left,
    numberstyle=\tiny\color{gray},
    tabsize=4,
    captionpos=t
    extendedchars=true,
    literate={а}{{\cyra}}1{б}{{\cyrb}}1{в}{{\cyrv}}1{г}{{\cyrg}}1{д}{{\cyrd}}1{е}{{\cyre}}1{ё}{{\cyryo}}1{ж}{{\cyrzh}}1{з}{{\cyrz}}1{и}{{\cyri}}1{й}{{\cyrishrt}}1{к}{{\cyrk}}1{л}{{\cyrl}}1{м}{{\cyrm}}1{н}{{\cyrn}}1{о}{{\cyro}}1{п}{{\cyrp}}1{р}{{\cyrr}}1{с}{{\cyrs}}1{т}{{\cyrt}}1{у}{{\cyru}}1{ф}{{\cyrf}}1{х}{{\cyrh}}1{ц}{{\cyrc}}1{ч}{{\cyrch}}1{ш}{{\cyrsh}}1{щ}{{\cyrshch}}1{ъ}{{\cyrhrdsn}}1{ы}{{\cyrery}}1{ь}{{\cyrsftsn}}1{э}{{\cyrerev}}1{ю}{{\cyryu}}1{я}{{\cyrya}}1{А}{{\CYRA}}1{Б}{{\CYRB}}1{В}{{\CYRV}}1{Г}{{\CYRG}}1{Д}{{\CYRD}}1{Е}{{\CYRE}}1{Ё}{{\CYRYO}}1{Ж}{{\CYRZH}}1{З}{{\CYRZ}}1{И}{{\CYRI}}1{Й}{{\CYRISHRT}}1{К}{{\CYRK}}1{Л}{{\CYRL}}1{М}{{\CYRM}}1{Н}{{\CYRN}}1{О}{{\CYRO}}1{П}{{\CYRP}}1{Р}{{\CYRR}}1{С}{{\CYRS}}1{Т}{{\CYRT}}1{У}{{\CYRU}}1{Ф}{{\CYRF}}1{Х}{{\CYRH}}1{Ц}{{\CYRC}}1{Ч}{{\CYRCH}}1{Ш}{{\CYRSH}}1{Щ}{{\CYRSHCH}}1{Ъ}{{\CYRHRDSN}}1{Ы}{{\CYRERY}}1{Ь}{{\CYRSFTSN}}1{Э}{{\CYREREV}}1{Ю}{{\CYRYU}}1{Я}{{\CYRYA}}1{№}{{\textnumero}}1
}

\titleformat{\section}{\normalfont\Large\bfseries}{\thesection.}{0.5em}{}
\titleformat{\subsection}{\normalfont\large\bfseries}{\thesubsection.}{0.5em}{}


\newcommand{\opencirclelocal}[2]{\filldraw[fill=white, draw=black] (axis cs:#1,#2) circle (1.5pt);}
\newcommand{\closedcirclelocal}[2]{\filldraw[fill=black, draw=black] (axis cs:#1,#2) circle (1.5pt);}

            
% Математические окружения
\theoremstyle{plain}
\newtheorem{theorem}{Теорема}[section]
\newtheorem{lemma}[theorem]{Лемма}
\newtheorem{state}{Утверждение}[section]

\theoremstyle{definition}
\newtheorem{definition}{Определение}[section]
\newtheorem{example}{Пример}[section]

\theoremstyle{remark}
\newtheorem{remark}{Замечание}[section]


\begin{document}

\begin{titlepage}
\begin{center}
    {\Large 
    ФЕДЕРАЛЬНОЕ ГОСУДАРСТВЕННОЕ БЮДЖЕТНОЕ \\
    ОБРАЗОВАТЕЛЬНОЕ \\
    УЧРЕЖДЕНИЕ ВЫСШЕГО ОБРАЗОВАНИЯ \\
    \vspace{0.5em} 
    «МОСКОВСКИЙ ГОСУДАРСТВЕННЫЙ УНИВЕРСИТЕТ \\
    имени М.В. ЛОМОНОСОВА»

    \vspace{3em} 

    МЕХАНИКО-МАТЕМАТИЧЕСКИЙ ФАКУЛЬТЕТ \\
    \vspace{1.5em} 
    КАФЕДРА ТЕОРИИ ВЕРОЯТНОСТЕЙ
    } 
\end{center}

\vspace{2em} 

\begin{center}
   \Large КУРСОВАЯ РАБОТА
\end{center} 

\begin{center}
    \Large 
    \textbf{ВЫХОДЯЩИЙ ПОТОК В СИСТЕМЕ МАССОВОГО ОБСЛУЖИВАНИЯ С БЕСКОНЕЧНЫМ ЧИСЛОМ КАНАЛОВ ПРИ ГРУППОВОМ ПОСТУПЛЕНИИ РАЗНОРОДНЫХ ЗАВИСИМЫХ ТРЕБОВАНИЙ}
\end{center}

\vfill


\noindent\hspace*{0.47\textwidth}
\begin{minipage}[t]{0.47\textwidth} 
    \raggedright 
    
    \large Выполнил: \\
    \large студент 408 группы \\
    \large Салимов Арсений Евгеньевич
    \vspace{2.5em} 

    \large Научный руководитель: \\ 
    \large профессор, доктор физ.-мат. наук \\ 
    \large Фалин Геннадий Иванович 
    \vspace{2.5em} 
    
\end{minipage}

\vspace*{\fill} 

\begin{center}
    Москва\\ \vspace{1em} 2026 г.
\end{center}

\end{titlepage}

\clearpage
\thispagestyle{empty}
\null
\newpage

\tableofcontents
\thispagestyle{empty}

\newpage
\mbox{} 
\thispagestyle{empty}

\newpage
\pagestyle{plain}
\setcounter{page}{1} 

\section*{Введение}
\addcontentsline{toc}{section}{Введение} % Чтобы попало в оглавление

Системы массового обслуживания с бесконечным числом приборов (типа $M/G/\infty$) являются фундаментальным инструментом моделирования процессов, где задержки обусловлены исключительно длительностью обслуживания, а не ожиданием в очереди. Благодаря своей аналитической гибкости, эти модели находят применение в самых разнообразных областях: от классических сетей связи до биологии и медицины.

Особый интерес для современных исследований представляют системы с \textit{групповым} поступлением заявок. Групповой трафик характерен для многих физических и информационных процессов. В недавних обзорах, например, в работе A.~Daw et al. (2022)~\cite{daw2022} выделяются ключевые направления использования таких моделей: архитектура микросервисов и облачные вычисления. Как показано в работе Y.~Gan и C.~Delimitrou~\cite{gan2018}, один пользовательский запрос к веб-приложению часто порождает группу взаимосвязанных подзадач, распределяемых по виртуальным серверам.

Критически важной характеристикой таких групп является их внутренняя структура. В реальных приложениях требования внутри одной группы часто являются \textit{гетерогенными} (разнородными) и имеют \textit{зависимые} времена обслуживания. Например, задачи, порожденные одним запросом (обращение к базе данных, логирование, вычисления), имеют принципиально разные распределения времени исполнения, но коррелированы по моменту поступления. Глубокий анализ таких систем с зависимым обслуживанием проводится в работе G.~Pang и W.~Whitt (2012)~\cite{pang2012}, где подчеркивается недостаточность простейших моделей с независимыми требованиями.

Исторически анализ систем с групповым поступлением проводился с использованием громоздкого аппарата производящих функций и дифференциальных уравнений. В классической работе Л.~М.~Абольникова (1968)~\cite{abolnikov1968} для нестационарной системы $M(t)/M/\infty$ использовались уравнения Колмогорова-Чепмена для производящих функций. Следующим шагом стало исследование B.~D.~Choi и K.~K.~Park (1992)~\cite{choi1992}, где была рассмотрена модель $M^k/M/\infty$ с гетерогенными требованиями. Авторам удалось найти распределение вектора числа требований в системе, однако они также использовали метод производящих функций, что привело к необходимости решения сложных систем дифференциальных уравнений в частных производных. Этот подход, будучи строгим, часто скрывал вероятностную суть процессов и был сложен для обобщения на случай произвольного распределения времени обслуживания.

Принципиально иной подход был предложен в работе Г.~И.~Фалина (1994)~\cite{falin1994}. Рассматривая модель $M^k/G/\infty$, где группы фиксированного размера $k$ поступают в соответствии с пуассоновским процессом, а требования внутри группы являются разнородными и зависимыми, автор доказал свойство декомпозиции. Было показано, что вектор числа требований различных типов в системе $N(t)$ может быть представлен в виде линейной комбинации независимых пуассоновских случайных величин. Это позволило свести сложную задачу анализа многомерного процесса к изучению характеристик отдельной группы, избегая громоздких выкладок.

Несмотря на то, что с момента публикации работы~\cite{falin1994} прошло несколько десятилетий, предложенный метод не теряет актуальности. В недавней работе A.~Daw, B.~Fralix и J.~Pender (2022)~\cite{daw2022} результаты Фалина были обобщены на случай нестационарных систем. Важно отметить, что в этой работе авторы исследовали не только число требований в системе, но и процесс ухода $D(t)$, получив выражение для его автоковариационной функции. Однако, в отличие от исходной постановки Фалина, авторы рассматривали упрощенный случай \textit{однородных} требований (неразличимых по типам).

При изучении динамики случайных процессов, таких как процесс обслуживания в СМО, помимо одномерных распределений, критически важную роль играют их корреляционные характеристики. Следуя классическому определению (\cite{bulinsky2005}, с. 53), ковариационная функция случайного процесса $X(t)$ определяется как:
\begin{equation}
    K(s, t) = \text{Cov}(X(s), X(t)) = \mathbb{E}\left[ (X(s) - \mathbb{E}X(s)) (X(t) - \mathbb{E}X(t)) \right].
\end{equation}

Целью настоящей курсовой работы является расширение анализа классической модели $M^k/G/\infty$, предложенной в работе Г.~И.~Фалина. В оригинальной статье основное внимание было сфокусировано на векторе числа требований в системе $N(t)$, в то время как выходящий процесс $D(t)$ (число обслуженных требований) остался нерассмотренным. Между тем, анализ выходного потока имеет ключевое значение для оценки нагрузки на последующие узлы сети. В данной работе будет показано, как вероятностный подход, основанный на свойстве расщепления пуассоновского потока, применяется для нахождения совместного распределения компонент вектора $D(t)$, анализа его ковариации с процессом $N(t)$. Также будет предложен метод для нахождения ковариационной функции процесса $D(t)$.

\newpage

\section{Постановка задачи}

Рассмотрим систему массового обслуживания $M^k/G/\infty$, где входящий поток заявок является пуассоновским процессом с параметром $\lambda$, и заявки поступают группами с фиксированным числом разнородных требований --- $k$, а число обслуживающих устройств бесконечно. Входящий поток групп будем обозначать через $A(t)$. Обозначим $\tau^{(n)}=(\tau_1^{(n)},...,\tau_k^{(n)})$ - случайный вектор времен обслуживания $n$-й поступившей группы. Мы будем предполагать, что вектора $\{\tau^{(n)}\}_{n=1}^{\infty}$ независимы, одинаково распределены и не зависят от входящего потока $A(t)$, однако их компоненты $\{\tau_i^{(n)}\}_{i=1}^k$ в общем случае являются зависимыми. Схематически нашу систему можно описать как показано на Рис. 1.

\begin{figure}[H]
    \centering
    \begin{tikzpicture}[
        >=latex, 
        scale=1,
        node distance=1.5cm,
        server/.style={circle, draw, thick, minimum size=1.2cm, fill=white},
        customer/.style={circle, draw, minimum size=0.5cm, inner sep=0pt, font=\small\bfseries},
        batch/.style={draw, dashed, rounded corners, inner sep=5pt, fill=gray!5}
    ]

    % --- 1. Входящий поток (Слева) ---
    \node (source) at (-4, 0) [align=center] {Входящий\\поток групп $A(t)$};
    
    % Стрелка с лямбдой наверху
    \draw[->, thick] (-2.3, 0) -- (-0.9, 0) node[midway, above] {$\lambda$};

    % --- 2. Группа (Центр) ---
    \node[batch] (batch) at (0, 0) {
        \begin{tikzpicture}
            \node[customer, fill=red!20] (c1) at (0, 1) {1};
            \node[customer, fill=green!20] (c2) at (0, 0) {2};
            \node at (0, -0.45) {$\vdots$};
            \node[customer, fill=blue!20] (ck) at (0, -1.1) {$k$};
        \end{tikzpicture}
    };
    \node[below=0.2cm] at (batch.south) {Группа ($n$-я)};

    % --- 3. Серверы (Справа, круги) ---
    \node[server] (s1) at (5, 2) {$\tau_1^{(n)}$};
    \node[server] (s2) at (5, 0) {$\tau_2^{(n)}$};
    \node at (5, -0.9) {$\vdots$};
    \node[server] (sk) at (5, -2) {$\tau_k^{(n)}$};

    % Стрелки распределения
    \draw[->, thick] (batch.east) -- (s1.west) node[midway, above, sloped, font=\footnotesize] {Тип 1};
    \draw[->, thick] (batch.east) -- (s2.west) node[midway, above, sloped, font=\footnotesize] {Тип 2};
    \draw[->, thick] (batch.east) -- (sk.west) node[midway, above, sloped, font=\footnotesize] {Тип $k$};

    % --- 4. Выход ---
    \draw[->, thick] (s1.east) -- (7, 2);
    \draw[->, thick] (s2.east) -- (7, 0);
    \draw[->, thick] (sk.east) -- (7, -2);
    
    % Объединяющая подпись справа (фигурная скобка)
    \draw[decorate, decoration={brace, amplitude=10pt}, thick, gray] (7.2, 2.2) -- (7.2, -2.2) 
        node[midway, right=15pt, align=left, black] {Выход:\\Вектор $D(t)$};

    \end{tikzpicture}
    \caption{Схематическое изображение системы обслуживания $M^k/G/\infty$ с разнородными требованиями.}
    \label{fig:schema_graph}
\end{figure}


Мы будем исследовать распределение случайного вектора $D(t) = (D_1(t),...,D_k(t))^T$, где $D_i(t)$  - число обслуженных требований $i$-го типа к моменту времени $t$. Будем считать, что в начальный момент времени $t=0$ в системе нет требований. Ввиду того, что вектора $\tau^{(n)}$ одинаково распределены, будет удобно ввести случайный вектор $\tau = (\tau_1,..., \tau_k)$, совпадающий по распределению с векторами $\tau^{(n)}$. 

Введем множество типов требований $S = \{1,..., k\}$. Моменты поступления групп обозначим $\alpha_1,\alpha_2,...$

\vspace{1em}
\begin{figure}[h!]
    \centering
    \begin{tikzpicture}[
        scale=1.2,
        every node/.style={font=\small},
        arrival/.style={circle, fill=black, inner sep=1.5pt},
        dep1/.style={circle, fill=red!80, inner sep=1.5pt},
        dep2/.style={circle, fill=blue!80, inner sep=1.5pt},
        arr_arrow/.style={-{Stealth[scale=1]}, dashed, thick}
    ]
        
        % Ось входящего потока (сверху)
        \draw[thick, ->] (0, 3) -- (10, 3) node[right] {$t$};
        \node[left, align=right] at (-0.2, 3) {Входящий поток\\групп $A(t)$};
        
        % Ось выходящего потока (снизу)
        \draw[thick, ->] (0, 0) -- (10, 0) node[right] {$t$};
        \node[left, align=right] at (-0.2, 0) {Выходящий поток\\требований $D(t)$};
        
        % --- Группа 1 (n=1) ---
        \coordinate (A1) at (1.5, 3);
        \node[arrival, label=above:{$\alpha_1$ (1-я группа)}] at (A1) {};
        
        % Выходы 1-й группы (допустим k=2)
        \coordinate (D1_1) at (3.5, 0); % Тип 1 выходит рано
        \coordinate (D1_2) at (6.5, 0); % Тип 2 выходит поздно
        
        \node[dep1, label=below:{$\alpha_1 + \tau_1^{(1)}$}] at (D1_1) {};
        \node[dep2, label=below:{$\alpha_1 + \tau_2^{(1)}$}] at (D1_2) {};
        
        \draw[arr_arrow, red!80] (A1) -- (D1_1) node[midway, left=2pt] {Тип 1};
        \draw[arr_arrow, blue!80] (A1) -- (D1_2);

        % --- Группа 2 (n=2) ---
        \coordinate (A2) at (4.5, 3);
        \node[arrival, label=above:{$\alpha_2$ (2-я группа)}] at (A2) {};
        
        % Выходы 2-й группы
        \coordinate (D2_1) at (5, 0); 
        \coordinate (D2_2) at (8.5, 0); 
        
        \node[dep1, label=below:{$\alpha_2 + \tau_1^{(2)}$}] at (D2_1) {};
        \node[dep2, label=below:{$\alpha_2 + \tau_2^{(2)}$}] at (D2_2) {};
        
        \draw[arr_arrow, red!80] (A2) -- (D2_1);
        \draw[arr_arrow, blue!80] (A2) -- (D2_2) node[midway, above right] {Тип 2};

    \end{tikzpicture}
    \caption{Входящий и выходящий поток во времени. Из-за случайного времени обслуживания $\tau_i^{(n)}$ требования из разных групп "перемешиваются" в выходящем потоке (например, требование 1-го типа из второй группы покинуло систему раньше, чем требование 2-го типа из первой).}
    \label{fig:timeline}
\end{figure}
\vspace{1em}

\section{Формулировка основного результата}
\begin{theorem}
$D(t) = B \cdot \delta(t)$, где:
\begin{enumerate}
    \item $\delta(t) = (\delta_y(t))_{y \in 2^S}$ - случайный вектор размера $2^k$ с независимыми компонентами.
    \item Каждая компонента $\delta_y(t)$ имеет распределение Пуассона с параметром $\lambda g_y(t)$, где
    \[
        g_y(t) = \int_0^t\mathbb{P}(\{\tau_i \le u\}_{i \in y} \cap \{\tau_i > u\}_{i \not\in y}) du
    \]
    \item 
    $B$ - матрица инцидентности размера $k \times 2^k$ с элементами
    \[
        b_{iy} =  \begin{cases}
            1,  & i \in y \\
            0, & i \not\in y
        \end{cases}
    \]
\end{enumerate}
\end{theorem}

\section{Доказательство теоремы}

Введем случайное подмножество $Y_n(t) \subseteq S$, характеризующее состояние $n$-ой группы в момент $t$:
\[
    Y_n(t) = \{i \in S : \alpha_n + \tau_i^{(n)} \le t\}
\]
Здесь событие $\{ i \in Y_n(t)\}$ означает, что требование типа $i$ из $n$-ой группы уже обслужено к моменту $t$. Наглядно процесс формирования множества $Y_n(t)$ в зависимости от положения момента времени относительно моментов завершения обслуживания представлен на Рис. 2.


\begin{figure}[H]
\centering
\begin{tikzpicture}[scale=1.2] % Масштаб для наглядности

% 1. Рисуем ось времени
\draw[-{Stealth}] (0,0) -- (11,0) node[below left] {$t$};

% 2. Отмечаем момент поступления требования alpha_n
\draw[thick] (2, -0.2) -- (2, 0.2);
\node[below=4pt] at (2, -0.2) {$\alpha_n$};

% 3. Отмечаем моменты завершения tau_i
% Момент завершения для i=1
\draw[thick] (4, -0.2) -- (4, 0.2);
\node[above=4pt] at (4, 0.2) {$\alpha_n + \tau_1^{(n)}$};

% Момент завершения для i=2
\draw[thick] (7, -0.2) -- (7, 0.2);
\node[above=4pt] at (7, 0.2) {$\alpha_n + \tau_2^{(n)}$};

% Момент завершения для i=3
\draw[thick] (9, -0.2) -- (9, 0.2);
\node[above=4pt] at (9, 0.2) {$\alpha_n + \tau_3^{(n)}$};

% 4. Показываем, как выглядит Y_n(t) в разные моменты времени

% --- Момент t_1 (до первого завершения) ---
\draw[dashed, red] (3, -1.2) -- (3, 1) node[above, black] {$t_1$};
\node[red, align=center] at (3, -1.5) {$Y_n(t_1) = \emptyset$};

% --- Момент t_2 (между первым и вторым завершением) ---
\draw[dashed, blue] (5.5, -1.2) -- (5.5, 1) node[above, black] {$t_2$};
\node[blue, align=center] at (5.5, -1.5) {$Y_n(t_2) = \{1\}$};

% --- Момент t_3 (после всех трех завершений) ---
\draw[dashed, green!60!black] (10, -1.2) -- (10, 1) node[above, black] {$t_3$};
\node[green!60!black, align=center] at (10, -1.5) {$Y_n(t_3) = \{1, 2, 3\}$};

\end{tikzpicture}
\caption{Визуализация множества $Y_n(t)$ в разные моменты времени $t$.}
\label{fig:time_axis}
\end{figure}

Теперь определим компоненты вектора $\delta(t)$:
\[
    \delta_y(t) = \sum_{n=1}^{A(t)}\mathbb{I}(Y_n(t)=y)
\]
$\delta_y(t)$ - это число групп, которые имеют исход $y$ к моменту $t$. Для демонстрации траекторий случайного процесса $\delta_y(t)$ и его связи с моментами обслуживания требований рассмотрим конкретный численный пример для системы с двумя типами требований ($k=2$) и тремя поступающими группами. Детальное описание сценария, таблица состояний $Y_n(t)$ и график траектории $\delta_y(t)$ приведены в \textbf{Приложении \ref{app:example}}.


Теперь выразим число обслуженных требований $i$-го типа в момент $t$ через $\delta_y(t)$:
\[
    D_i(t) = \sum_{n=1}^{A(t)}\mathbb{I}(i \in Y_n(t)) = \sum_{y \in 2^S}\sum_{n=1}^{A(t)}\mathbb{I}(i \in Y_n(t), Y_n(t) = y) = \sum_{y\in 2^S: i \in y}\sum_{n=1}^{A(t)}\mathbb{I}(Y_n(t)=y) = 
\]
\[
    = \sum_{y\in 2^S: i \in y}\delta_y(t)
\]

Отсюда следует, что $D(t) = B \cdot \delta(t)$. Теперь найдем распределение компонент вектора $\delta(t)$.

Обозначим $\pi(\overline{n}, t) = \mathbb{P}(\delta(t) = \overline{n})$, где $\overline{n} = (n_y)_{y \in 2^S}$, $n_y \ge 0$, $|\overline{n}| = \sum_{y \in 2^S}n_y$. 

Заметим, что 
\[
    \sum_{y \in 2^S} \delta_y(t) = \sum_{y \in 2^S}\sum_{n=1}^{A(t)}\mathbb{I}(Y_n(t) = y) = \sum_{n=1}^{A(t)}\sum_{y \in 2^S}\mathbb{I}(Y_n(t) = y) = \sum_{n=1}^{A(t)} 1 = A(t) 
\]
так как во внутренней сумме лишь один индикатор равен 1, а все остальные равны 0.

Тогда событие $\{\delta(t) = \overline{n}\}$ вложено в событие $\{A(t) = |\overline{n}|\}$ и, применив формулу полной вероятности, можно записать, что
\[
    \pi(\overline{n}, t) = \mathbb{P}(\delta(t) = \overline{n}, A(t) = |\overline{n}|) = \mathbb{P}(\delta(t) = \overline{n}| A(t) = |\overline{n}|) \cdot \mathbb{P}(A(t) = |\overline{n}|)
\]

Поскольку $A(t) \sim \Pi(\lambda t)$, то $\mathbb{P}(A(t) = |\overline{n}|) = \frac{(\lambda t)^{|\overline{n}|}}{|\overline{n}|!}e^{-\lambda t}$

Также известно(\cite{falinrisk}, раздел 3.4), что если в пуассоновском процессе за интервал $(0,t]$ поступило $m$ групп, то моменты их поступления распределены как $m$ порядковых статистик независимых равномерно распределенных на $(0,t]$ случайных величин $U_1,...,U_m$.

Определим случайное множество $Z = \{i \in S: U + \tau_i \le t\}$, где $U$ равномерно распределена на $(0, t]$. Это множество является исходом произвольной группы прибывшей на интервале $(0,t]$. Обозначим $\pi_y(t) = \mathbb{P}(Z = y)$ и найдем эти вероятности. Для группы, поступившей на интервале $(0, t]$, исход $y$ наступает, если $U + \tau_i \le t$ для $i \in y$ и $U  + \tau_i > t$ для $i \not\in y$. По формуле полной вероятности для непрерывного случая, получим, что 
\[
    \pi_y(t) = \int_0^t\mathbb{P}(\{\tau_i \le t - u)\}_{i \in y} \cap \{\tau_i > t - u\}_{i \not\in y})\frac{1}{t}du
\]
После замены $t - u = \xi$, получим
\[
    \pi_y(t) = \frac{1}{t}\int_0^t\mathbb{P}(\{\tau_i \le \xi)\}_{i \in y} \cap \{\tau_i > \xi\}_{i \not\in y})d\xi = \frac{1}{t}g_y(t)
\]

Поскольку все группы между собой независимы, то условное распределение $\delta(t)$ является полиномиальным с вероятностями $\pi_y(t)$ и 
\[
    \mathbb{P}(\delta(t) = \overline{n}| A(t) = |\overline{n}|) = \frac{|\overline{n}|}{\prod_y n_y!}\prod_y(\pi_y(t))^{n_y}
\]

Теперь заметим, что $\sum_{y\in 2^S}g_y(t) = t$:
\[
    \sum_{y\in 2^S}g_y(t) = \int_0^t \sum_{y \in 2^S}\mathbb{P}(\{\tau_i \le u\}_{i \in y} \cap \{\tau_i > u\}_{i \not\in y}\} du
\]
События при разных $y$ не пересекаются и их объединение дает все вероятностное пространство, поэтому вся сумма превратится в 1.

Тогда

\[
    \pi(\overline{n}, t) = \frac{|\overline{n}|!}{\prod_y n_y!}\prod_y(\pi_y(t))^{n_y} \cdot \frac{(\lambda t)^{|\overline{n}|}}{|\overline{n}|!}e^{-\lambda t} = \prod_y \frac{(g_y(t))^{n_y}}{n_y!t^{n_y}} \cdot \lambda^{|\overline{n}|} t^{|\overline{n}|} e^{-\lambda t} = 
\]
\[
    = \prod_y \frac{(\lambda g_y(t))^{n_y}}{n_y!}e^{-\lambda g_y(t)}
\]
Теперь зафиксируем множество $y^*$ и просуммируем по всем остальным множествам $y$. В результате получим, что
\[
    \delta_{y^*}(t) \sim \Pi(\lambda g_{y^*}(t))
\]
и компоненты вектора $\delta(t)$ независимы.

\section{Ковариация процессов $N(t)$ и $D(t)$}

Теперь найдем ковариацию векторов $N(t)$ и $D(t)$. Для этого будет удобно выразить компоненты вектора $N(t)$ через сумму независимых пуассоновских случайных величин $\delta_y(t)$. Вспомним определение случайного подмножества $Y_n(t)$:
\[
    Y_n(t) = \{i \in S: \alpha_n + \tau_i^{(n)} \le t\}
\]
Заметим, что событие $\{i \not \in Y_n(t)\}$ само по себе означает лишь то, что обслуживание требования $i$ из $n$-й группы еще не завершено (либо группа еще не поступила). Однако, поскольку суммирование ведется по $n$ от $1$ до $A(t)$, (то есть только по уже поступившим группам), условие $\{i \not \in Y_n(t)\}$ однозначно гарантирует, что требование находится в системе в момент времени $t$. Тогда 
\[
    N_i(t) = \sum_{n=1}^{A(t)} \mathbb{I}(i \not \in Y_n(t)) = \sum_{y \in 2^S} \sum_{n=1}^{A(t)}\mathbb{I}(i \not \in Y_n(t), Y_n(t) = y) = \sum_{y \in 2^S, i \not \in y}\sum_{n = 1}^{A(t)}\mathbb{I}(Y_n(t) = y) =
\]
\[
    =  \sum_{y \in 2^S, i \not \in y} \delta_y(t)
\]
Таким образом, мы получили схожее выражение для $N_i(t)$. В матричной форме это выражение для вектора $N(t)$ будет иметь следующий вид:
\[
    N(t) = (J - B)\cdot\delta(t),
\]
где $J$ - матрица размера $k \times 2^k$ полностью состоящая из единиц.

Теперь перейдем к нахождению ковариации:
\[
    \mathrm{Cov}(N_i(t), D_j(t)) = \mathrm{Cov}\bigg(\sum_{y \in 2^S: i \in y}\delta_y(t), \sum_{z \in 2^S : j \not \in z}\delta_z(t)\bigg) = \sum_{y \in 2^S: i \in y}\sum_{z \in 2^S : j \not \in z}\mathrm{Cov}\bigg(\delta_y(t), \delta_z(t)\bigg)
\]
Поскольку величины $\delta_y(t)$ попарно независимы, то их ковариация при $y \not = z$ равна 0. При $y = z$ получим дисперсию $\delta_y(t)$, равную $\lambda g_y(t)$, так как $\delta_y(t) \sim \Pi(\lambda g_y(t))$. В итоге получим:
\begin{equation}
    \mathrm{Cov}(N_i(t), D_j(t)) = \lambda \sum_{\substack{y \in 2^S \\ i \in y \\ j \not \in y}}g_y(t) \label{eq:cov}
\end{equation}
Теперь рассмотрим два случая:
\begin{enumerate}
    \item $i = j$ \\
    В таком случае индекс $i$ должен одновременно принадлежать и не принадлежать множествам, по которым идет суммирование. Значит вся сумма равна 0, то есть $\mathrm{Cov}(N_i(t), D_i(t)) = 0$.

    Более того, в этом случае величины $N_i(t)$ и $D_i(t)$ являются независимыми. Действительно, 
    \[
        N_i(t) = \sum_{y \in 2^S, i \not \in y} \delta_y(t)
    \]
    \[
        D_i(t) = \sum_{y \in 2^S, i \in y} \delta_y(t)
    \]
    Множества индексов суммирования не пересекаются и все величины $\delta_y(t)$ являются независимыми. По лемме о группировке $N_i(t)$ и $D_i(t)$ независимы.
    \item $i \not = j$ \\
    Напомним определение функций $g_y(t)$:
    \[
        g_y(t) = \int_0^t \mathbb{P}(\{\tau_m \le u\}_{m \in y} \cap \{\tau_m > u\}_{m \not \in y})du \ge 0
    \]
    Если времена обслуживания таковы, что какая-то из этих вероятностей строго больше $0$, то $\mathrm{Cov}(N_i(t), D_j(t))$ > 0, а значит величины зависимы.
\end{enumerate}

Поскольку $\mathrm{Cov}(Ax, By) = A \cdot\mathrm{Cov}(x, y) \cdot B^T$, то ковариационная матрица $C(t)$ будет иметь следующее представление:
\[
    C(t) = \mathrm{Cov}(N(t), D(t)) = \mathrm{Cov}((J - B)\cdot\delta(t), B\cdot \delta(t)) = (J - B) \cdot \mathrm{Cov}(\delta(t), \delta(t)) \cdot B^T = 
\]
\[
    = (J - B) \cdot \Sigma_\delta(t) \cdot B^T,
\]
где $\Sigma_\delta(t)$ - диагональная матрица размера $2^k \times 2^k$, где на диагонали стоят дисперсии $\lambda g_y(t)$.

Таким образом, мы можем охарактеризовать зависимость между векторами $N(t)$ и $D(t)$ через их ковариационную матрицу $C(t)$, причем
\begin{enumerate}
    \item диагональные элементы $c_{ii}$ равны 0, что отражает независимость числа обслуживающихся и обслуженных требований одного типа
    \item внедиагональные элементы $c_{ij}$ при $i \not = j$ строго положительны(при выполнении условия $g_y(t) > 0$), что подтверждает зависимость между числом обслуживаемых и обслуженных требований разных типов. 
\end{enumerate}

\begin{example}
Проиллюстрируем вид ковариационной матрицы при $k=2$. Пусть случайный вектор времен обслуживания $\tau = (\tau_1, \tau_2)$ удовлетворяет следующим условиям:
\begin{enumerate}
    \item Маргинальные распределения $\tau_1$ и $\tau_2$ являются экспоненциальными с параметром $\mu = 1$.
    \item Компоненты $\tau_1$ и $\tau_2$ зависимы, а их совместное распределение принадлежит классу Фарли-Гумбела-Моргенштерна (FGM):
\end{enumerate}
\[
    F(t_1, t_2) = F_1(t_1) F_2(t_2) \left[1 + \alpha \big(1-F_1(t_1)\big)\big(1-F_2(t_2)\big)\right],
\]
где $F_i(t_i) = 1 - e^{-t_i}$ при $t_i \ge 0$, а параметр $\alpha$ лежит в интервале $[-1, 1]$. 
В явном виде формула записывается так:
\[
    F(t_1, t_2) = (1- e^{-t_1})(1 - e^{-t_2}) \left(1 + \alpha e^{-t_1 - t_2}\right).
\]

\begin{remark}
Ограничение $|\alpha| \le 1$ является строгим: при выходе за эти границы совместная плотность вероятности может принимать отрицательные значения, что недопустимо.
\end{remark}

Как показано в работе \cite{Gumbel}, для такого распределения $\mathrm{Cov}(\tau_1, \tau_2) = \alpha/4$. Из этого факта и структуры функции распределения следует, что времена обслуживания независимы тогда и только тогда, когда $\alpha = 0$.

Чтобы наглядно продемонстрировать влияние параметра $\alpha$ на структуру зависимости времён обслуживания внутри одной группы, рассмотрим совместную плотность вероятности $f(t_1, t_2) = \frac{\partial^2}{\partial t_1 \partial t_2} F(t_1, t_2)$. Визуализация линий уровня (изолиний) данной плотности представлена на Рис. \ref{fig:fgm_contours} (исходный код для генерации графиков приведен в Приложении \ref{app:python_code}).

\begin{figure}[h]
    \centering
    \includegraphics[width=\textwidth]{Images/fgm_density_contours.pdf} 
    \caption{Линии уровня совместной плотности вероятности времён обслуживания $f(t_1, t_2)$ для распределения FGM при различных значениях параметра зависимости $\alpha$.}
    \label{fig:fgm_contours}
\end{figure}

Анализируя представленные графики, можно выделить следующие режимы работы системы обслуживания:
\begin{itemize}
    \item При $\alpha = 0$ (центральный график) времена обслуживания независимы. Совместная плотность имеет вид $f(t_1, t_2) = e^{-(t_1+t_2)}$, поэтому линии постоянной вероятности представляют собой прямые отрезки.
    
    \item При $\alpha > 0$ (правый график, $\alpha = 1$) наблюдается положительная корреляция. Изолинии плотности становятся выпуклыми и «вытягиваются» вдоль диагонали $t_1 = t_2$. В контексте систем массового обслуживания это означает, что требования внутри одной группы имеют схожую сложность: если обслуживание первого требования затягивается, с высокой вероятностью затянется и обслуживание второго.
    
    \item При $\alpha < 0$ (левый график, $\alpha = -1$) корреляция отрицательна. Изолинии принимают вогнутую (L-образную) форму, прижимаясь к осям координат. Это моделирует ситуацию конкурирующих требований: длительное обслуживание одного требования из группы обуславливает быстрое завершение второго.
\end{itemize}

Перейдем к вычислению ковариационной матрицы. При $k = 2$ множество подмножеств индексов $2^S$ состоит из элементов $\emptyset, \{1\}, \{2\}, \{1, 2\}$. В формуле \eqref{eq:cov} слагаемые отличны от нуля только при выполнении условий $i \in y$ и $j \notin y$.
\begin{itemize}
    \item Для $y = \emptyset$ условие $i \in y$ всегда ложно.
    \item Для $y = \{1, 2\}$ условие $j \notin y$ ложно для любых $i, j \in \{1, 2\}$.
\end{itemize}
Следовательно, достаточно вычислить функцию $g_y(t)$ только для одноэлементных множеств $y = \{1\}$ и $y = \{2\}$.

\begin{enumerate}
    \item Найдем $g_{\{1\}}(t)$:
    \[
        g_{\{1\}}(t) = \int_0^t \mathbb{P}(\tau_1 \le u, \tau_2 > u) \, du.
    \]
    Вероятность под интегралом можно выразить следующим образом:
    \[
        \mathbb{P}(\tau_1 \le u, \tau_2 > u) = \mathbb{P}(\tau_2 > u) - \mathbb{P}(\tau_1 > u, \tau_2 > u).
    \]
    Подставив явный вид распределения FGM, получим:
    \[
        \mathbb{P}(\tau_1 \le u, \tau_2 > u) = e^{-u} - \left[ (\alpha + 1)e^{-2u} - 2\alpha e^{-3u} + \alpha e^{-4u} \right].
    \]
    Интегрируя это выражение по $u$ от 0 до $t$, получаем окончательную формулу для $g_{\{1\}}(t)$:
    \[
        g_{\{1\}}(t) = (1 - e^{-t}) - \frac{\alpha + 1}{2}(1-e^{-2t}) + \frac{2\alpha}{3}(1-e^{-3t}) - \frac{\alpha}{4}(1 - e^{-4t}).
    \]

    \item В силу симметрии распределения относительно $\tau_1$ и $\tau_2$, выражение для $g_{\{2\}}(t)$ совпадает с $g_{\{1\}}(t)$.
\end{enumerate}

Перейдем к формированию ковариационной матрицы $C(t) = \mathrm{Cov}(N(t), D(t))$.
Напомним, что диагональные элементы матрицы (ковариации процессов одного типа) в данной модели равны нулю. Вычислим внедиагональные элементы:
\[
    \mathrm{Cov}(N_1(t), D_2(t)) = \lambda \sum_{\substack{y \in 2^S \\ 1 \in y \\ 2 \notin y}} g_y(t) = \lambda g_{\{1\}}(t).
\]
Аналогично, в силу симметрии:
\[
    \mathrm{Cov}(N_2(t), D_1(t)) = \lambda \sum_{\substack{y \in 2^S \\ 2 \in y \\ 1 \notin y}} g_y(t) = \lambda g_{\{2\}}(t) = \lambda g_{\{1\}}(t).
\]
Таким образом, ковариационная матрица является симметричной и имеет вид:
\[
C(t) = 
\begin{pmatrix}
    0 & \lambda g_{\{1\}}(t) \\
    \lambda g_{\{1\}}(t) & 0
\end{pmatrix}.
\]


\begin{remark}
    Важно отметить, что даже при независимых временах обслуживания (случай $\alpha = 0$) процессы $N(t)$ и $D(t)$ остаются зависимыми. Подставив $\alpha = 0$ в выражение для $g_{\{1\}}(t)$, получаем:
    \[
        g_{\{1\}}(t)\bigg|_{\alpha=0} = 1 - e^{-t} - \frac{1}{2}(1-e^{-2t}) = \frac{1}{2}(1-e^{-t})^2 > 0.
    \]
    Это означает наличие корреляции между выходящим потоком и потоком обслуживанных требований даже при отсутствии корреляции во времени обслуживания.
\end{remark}

Заметим, что при $t \to \infty$ выражение для $g_{\{1\}}(t)$ имеет предел, равный
\[
    \lim_{t \to \infty}g_{\{1\}}(t) = 1 - \frac{\alpha + 1}{2} + \frac{2\alpha}{3}- \frac{\alpha}{4} = \frac{1}{2} - \frac{\alpha}{12}
\]

Полученный предел имеет ясный вероятностный и физический смысл. Поскольку случайная величина $\delta_{\{1\}}(t)$ (число групп, в которых первое требование обслужено, а второе --- еще нет) имеет распределение Пуассона с параметром $\lambda g_{\{1\}}(t)$, существование конечного предела при $t \to \infty$ доказывает, что процесс числа «частично обслуженных» групп имеет стационарное распределение. В стационарном режиме среднее число таких групп в системе постоянно и равно $\lambda \left(\frac{1}{2} - \frac{\alpha}{12}\right)$.

Из этой формулы наглядно видно влияние зависимости времён обслуживания:
\begin{itemize}
    \item При \textbf{положительной корреляции} ($\alpha > 0$) времена обслуживания $\tau_1$ и $\tau_2$ склонны принимать близкие значения. Следовательно, интервал времени, в течение которого группа находится в состоянии $\{1\}$, сокращается, и среднее число таких групп в системе падает (достигая минимума $\frac{5}{12}\lambda$ при $\alpha = 1$).
    
    \item При \textbf{отрицательной корреляции} ($\alpha < 0$) одно требование обслуживается быстро, а другое --- долго. Время, когда одно требование уже покинуло систему, а второе еще находится в ней, увеличивается, что приводит к увеличению числа таких групп (достигая максимума $\frac{7}{12}\lambda$ при $\alpha = -1$).
    
    \item Случай независимых времён обслуживания ($\alpha = 0$) дает промежуточное значение $\frac{1}{2}\lambda$.
\end{itemize}

Динамика выхода функции $g_{\{1\}}(t)$ на стационарные асимптоты для различных значений параметра $\alpha$ проиллюстрирована на Рис.~\ref{fig:g1_plot}(соответствующий программный код приведен в Приложении \ref{app:python_code}).

\vspace{1em}
\begin{figure}[H]
    \centering
    \includegraphics[width=0.85\textwidth]{Images/g1_plot.pdf}
    \caption{Зависимость функции $g_{\{1\}}(t)$ от времени $t$ при различных значениях параметра $\alpha$. Пунктирными линиями показаны горизонтальные асимптоты $\frac{1}{2} - \frac{\alpha}{12}$, к которым стремятся вероятности при $t \to \infty$. График иллюстрирует процесс выхода числа «частично обслуженных» групп на стационарный режим.}
    \label{fig:g1_plot}
\end{figure}
\vspace{1em}
\end{example}

\section{Ковариационная функция процесса $D(t)$}

Формально выходящий процесс $D(t)$ может быть представлен как сумма случайного числа независимых случайных векторов следующим образом:
\[
    D(t) = \sum_{n=1}^{A(t)}Z_n(t),
\]
где $Z_n(t)$ - вектор индикаторов обслуживания для $n-$й группы, то есть $(Z_n(t))_i = \mathbb{I}(\alpha_n + \tau_i^{(n)} \le t)$. То есть мы получили представление общего вектора обслуженных требований через сумму векторов, каждый из которых относится только к $n-$й поступившей группе, а поскольку группы между собой не зависимы, то таким образом можно было бы упростить вычисление автоковариации. Однако такой способ ведет к некоторым техническим трудностям:
\[
    \mathrm{Cov}(D(t_1), D(t_2)) = \mathrm{Cov}\bigg(\sum_{n=1}^{A(t_1)}Z_n(t_1), \sum_{n=1}^{A(t_2)}Z_m(t_2)\bigg)
\]
Трудность заключается в том, что для случайных сумм мы не можем пользоваться свойством линейности ковариации и вынести обе суммы за ковариацию. Имеет место следующий контрпример.
\begin{example}
    Пусть $N \sim \Pi(m)$ и $S = \sum_{n=1}^N 1$, то есть $S = N$. Теперь подсчитаем дисперсию $S$, то есть $\mathrm{Cov}(S, S)$:
    \[
        \mathrm{Var}(S) = \mathrm{Var}(N) = m > 0
    \]
    А если попробовать вынести сумму за ковариацию, то получится следующее:
    \[
        \mathrm{Var}(S) = \mathrm{Cov}\bigg(\sum_{n=1}^N 1, \sum_{m=1}^N 1 \bigg) = \sum_{n=1}^N\sum_{m=1}^N\mathrm{Cov}(1, 1) = 0
    \]
    Что ведет к противоречию с $m > 0$.
\end{example}

Здесь можно было бы применить формулу полной ковариации
\[
    \mathrm{Cov}(X, Y) = \mathbb{E}(\mathrm{Cov}(X, Y | A)) + \mathrm{Cov}(\mathbb{E}(X | A), \mathbb{E}(Y | A)),
\]
однако последующие выкладки получились бы очень громоздкими. 

Для нахождения ковариационной функции выходящего процесса будет использован совершенно иной подход. Пусть $t_1 \le t_2$. Разобьем отрезок $[0, t_2]$ на маленькие промежутки длины $\Delta u$. Из свойств процесса Пуассона с интенсивностью $\lambda$ за промежуток длины $\Delta u$ при $\Delta u \to 0$:
\begin{enumerate}
    \item Вероятность прихода ровно 1 группы равна $\lambda \Delta u + o(\Delta u)$
    \item С вероятностью $1 - \lambda \Delta u + o(\Delta u)$ новая группа не поступит
    \item Вероятность прихода более 1 группы равна $o(\Delta u)$.
\end{enumerate}

Обозначим через $\Delta D(t, u)$ - вклад в процесс $D(t)$ групп, которые поступают на промежутке $[u, u + \Delta u)$. Тогда 
\[
    D(t) = \sum_n\Delta D(t, u_n), \text{ где } \{u_n\} \text{ - разбиение отрезка } [0, t_2] 
\]
Поскольку вероятность поступления более одной группы на таком интервале равна $o(\Delta u)$, то будем считать, то группа либо не поступает вовсе, либо поступает ровно одна группа. Тогда компоненты вектора $\Delta D(t, u)$ можно представить следующим образом:
\[
    \Delta D_i(t, u) = \eta(u) \cdot \mathbb{I}(\text{требование типа } i \text{ обслужено к моменту } t),
\]
где $\eta(u)$ принимает значения $0$ и $1$ с вероятностями $1 - \lambda \Delta u + o(\Delta u)$ и $\lambda \Delta u + o(\Delta u)$ соответственно. Тогда 
\[
    \mathrm{Cov}(D(t_1), D(t_2)) = \mathrm{Cov}\bigg(\sum_n\Delta D(t_1, u_n), \sum_m\Delta D(t_2, u_m)\bigg) = 
\]
\[
    = \sum_n \sum_m\mathrm{Cov}(\Delta D(t_1, u_n), \Delta D(t_2, u_m))
\]
Поскольку входящий поток является пуассоновским, то поступления в разные интервалы $u_n \not = u_m$ являются независимыми и слагаемые при $n \not = m$ будут равны 0.
\[
    \mathrm{Cov}(D(t_1), D(t_2)) = \sum_n\mathrm{Cov}(\Delta D(t_1, u_n), \Delta D(t_2, u_n))
\]
Заметим, что если какой-то $u_n > t_1$, то слагаемое в сумме будет равно 0, так как $\Delta D(t, u)$ отличен от нулевого вектора только при $u \le t$. Соответственно, суммирование будет идти только по тем точкам $u_n$, которые на числовой оси находятся левее, чем $t_1$. 

Теперь вычислим каждое слагаемое по отдельности. Обозначим $p = \lambda \Delta u$ и пусть $I_i, I_j$ - индикаторы обслуживания требований типа $i, j$, если группа пришла:
\[
    \mathrm{Cov}(\Delta D_i(t_1, u), \Delta D_j(t_2, u)) = \mathrm{Cov}(\eta(u)I_i, \eta(u) I_j) = \mathbb{E}(\eta(u)I_i \cdot \eta(u) I_j) - \mathbb{E}(\eta(u)I_i)\cdot \mathbb{E}(\eta(u)I_j)
\]

Так как $\eta(u)^2 = \eta(u)$ и она не зависит от $I_i, I_j$, получаем:
\[
    p\mathbb{E}(I_iI_j) - p^2\mathbb{E}I_i\mathbb{E}I_j + o(\Delta u) = \lambda \Delta u\mathbb{E}(I_iI_j) + o(\Delta u)
\]
Случайная величина $I_iI_j$ является индикатором того, что требования типов $i, j$, поступившие на интервале $[u, u + \Delta u)$ будут обслужены к моменту $t_1$ и $t_2$. Переходя к пределу по $\Delta u \to 0$ получим следующий интеграл:
\[
    \mathrm{Cov}(D_i(t_1), D_j(t_2)) = \lambda \int_0^{t_1}\mathbb{P}(u + \tau_i \le t_1, u + \tau_j \le t_2) du
\]
При $i = j$ имеем: 
\[
    \mathrm{Cov}(D_i(t_1), D_i(t_2)) = \lambda \int_0^{t_1}\mathbb{P}(u + \tau_i \le t_1)du
\]

\newpage

\section[Явный вид ковариационной матрицы процесса $D(t)$ \protect\newline для FGM-распределения]{Явный вид ковариационной матрицы процесса $D(t)$ для FGM-распределения}

В предыдущем разделе было получено общее интегральное представление для ковариации процесса $D(t)$. Чтобы продемонстрировать работоспособность данного результата, вычислим ковариационную матрицу $C(t_1, t_2) = \text{Cov}(D(t_1), D(t_2))$ явно для двумерного случая ($k=2$). Предположим, что времена обслуживания $\tau_1$ и $\tau_2$ имеют совместное распределение Фарли-Гумбела-Моргенштерна с экспоненциальными маргиналами (параметр $\mu=1$), как было описано в Примере 4.1:
\begin{equation}
    F(x, y) = \mathbb{P}(\tau_1 \le x, \tau_2 \le y) = (1 - e^{-x})(1 - e^{-y}) \left[ 1 + \alpha e^{-x-y} \right], \quad x, y \ge 0.
\end{equation}

Ковариационная матрица процесса $D(t)$ размера $2 \times 2$ состоит из элементов $c_{ij}(t_1, t_2) = \text{Cov}(D_i(t_1), D_j(t_2))$. Без ограничения общности будем считать, что $t_1 \le t_2$. Вычислим отдельно ее диагональные и внедиагональные элементы.

\subsection{Диагональные элементы матрицы ковариации}
Для диагональных элементов $c_{ii}(t_1, t_2)$, где $i \in \{1, 2\}$, формула из конца раздела 5 принимает вид:
\begin{equation}
    c_{ii}(t_1, t_2) = \lambda \int_0^{t_1} \mathbb{P}(u + \tau_i \le t_1) du.
\end{equation}

Поскольку маргинальные распределения $\tau_i$ являются экспоненциальными с параметром $\mu=1$, вероятность под интегралом равна $F_i(t_1 - u) = 1 - e^{-(t_1-u)}$. Вычислим интеграл:
\[
    c_{ii}(t_1, t_2) = \lambda \int_0^{t_1} \left( 1 - e^{-t_1+u} \right) du = \lambda \left[ u - e^{-t_1}e^u \right]_0^{t_1} = 
\]
\begin{equation}
    = \lambda \left( (t_1 - 1) - (0 - e^{-t_1}) \right) = \lambda (t_1 - 1 + e^{-t_1}).
\end{equation}

Заметим, что в общем случае для произвольных $t_1$ и $t_2$ результат зависит только от минимального из времен, то есть $c_{11}(t_1, t_2) = c_{22}(t_1, t_2) = \lambda (\min(t_1, t_2) - 1 + e^{-\min(t_1, t_2)})$.

\subsection{Внедиагональные элементы матрицы ковариации}
Теперь вычислим элемент $c_{12}(t_1, t_2) = \text{Cov}(D_1(t_1), D_2(t_2))$, описывающий ковариацию между процессами обслуживания требований первого и второго типов. Согласно формуле из раздела 5:
\begin{equation}
    c_{12}(t_1, t_2) = \lambda \int_0^{t_1} F(t_1-u, t_2-u) du.
\end{equation}

Подставим совместное распределение FGM в подынтегральное выражение и раскроем скобки:
\begin{align*}
    F(t_1-u, t_2-u) &= \left(1 - e^{-t_1+u}\right)\left(1 - e^{-t_2+u}\right) \left[ 1 + \alpha e^{-t_1-t_2+2u} \right] = \\
    &= \left( 1 - e^{-t_1}e^u - e^{-t_2}e^u + e^{-t_1-t_2}e^{2u} \right) + \\
    &\quad + \alpha \left( e^{-t_1-t_2}e^{2u} - e^{-2t_1-t_2}e^{3u} - e^{-t_1-2t_2}e^{3u} + e^{-2t_1-2t_2}e^{4u} \right).
\end{align*}

Интегрирование данного полинома от экспонент по $u$ на отрезке $[0, t_1]$ сводится к использованию табличных интегралов $\int_0^{t_1} e^{ku} du = \frac{1}{k}(e^{kt_1} - 1)$. Для наглядности представим результат интегрирования в виде суммы двух компонент: $c_{12}(t_1, t_2) = \lambda \left( I_0 + \alpha I_\alpha \right)$, где $I_0$ — базовая часть, не зависящая от $\alpha$, а $I_\alpha$ — поправка, обусловленная зависимостью времен обслуживания.

Вычислим базовую часть $I_0$:
\begin{align*}
    I_0 &= \int_0^{t_1} \left( 1 - e^{-t_1}e^u - e^{-t_2}e^u + e^{-t_1-t_2}e^{2u} \right) du = \\
    &= t_1 - e^{-t_1}(e^{t_1}-1) - e^{-t_2}(e^{t_1}-1) + \frac{1}{2}e^{-t_1-t_2}(e^{2t_1}-1) = \\
    &= t_1 - 1 + e^{-t_1} - e^{t_1-t_2} + e^{-t_2} + \frac{1}{2}e^{t_1-t_2} - \frac{1}{2}e^{-t_1-t_2} = \\
    &= t_1 - 1 + e^{-t_1} + e^{-t_2} - \frac{1}{2}e^{t_1-t_2} - \frac{1}{2}e^{-t_1-t_2}.
\end{align*}

Теперь вычислим компоненту $I_\alpha$:
\begin{align*}
    I_\alpha &= \int_0^{t_1} \left( e^{-t_1-t_2}e^{2u} - e^{-2t_1-t_2}e^{3u} - e^{-t_1-2t_2}e^{3u} + e^{-2t_1-2t_2}e^{4u} \right) du = \\
    &= \frac{1}{2}e^{-t_1-t_2}(e^{2t_1}-1) - \frac{1}{3}e^{-2t_1-t_2}(e^{3t_1}-1) - \frac{1}{3}e^{-t_1-2t_2}(e^{3t_1}-1) + \frac{1}{4}e^{-2t_1-2t_2}(e^{4t_1}-1).
\end{align*}
Раскрывая скобки и группируя слагаемые с одинаковыми показателями степени, получаем:
\begin{align*}
    I_\alpha &= \left( \frac{1}{2} - \frac{1}{3} \right)e^{t_1-t_2} + \left( \frac{1}{4} - \frac{1}{3} \right)e^{2t_1-2t_2} - \\
    &\quad - \frac{1}{2}e^{-t_1-t_2} + \frac{1}{3}e^{-2t_1-t_2} + \frac{1}{3}e^{-t_1-2t_2} - \frac{1}{4}e^{-2t_1-2t_2} = \\
    &= \frac{1}{6}e^{t_1-t_2} - \frac{1}{12}e^{2(t_1-t_2)} - \frac{1}{2}e^{-t_1-t_2} + \frac{1}{3}e^{-2t_1-t_2} + \frac{1}{3}e^{-t_1-2t_2} - \frac{1}{4}e^{-2t_1-2t_2}.
\end{align*}

Таким образом, внедиагональный элемент ковариационной матрицы равен:
\begin{align}
    c_{12}(t_1, t_2) &= \lambda \Bigg[ t_1 - 1 + e^{-t_1} + e^{-t_2} - \frac{1}{2}e^{t_1-t_2} - \frac{1}{2}e^{-t_1-t_2} + \nonumber \\
    &\quad + \alpha \left( \frac{1}{6}e^{t_1-t_2} - \frac{1}{12}e^{2(t_1-t_2)} - \frac{1}{2}e^{-t_1-t_2} + \frac{1}{3}e^{-2t_1-t_2} + \frac{1}{3}e^{-t_1-2t_2} - \frac{1}{4}e^{-2t_1-2t_2} \right) \Bigg]. \label{eq:cov_matrix_offdiag}
\end{align}

\subsection{Итоговый вид матрицы и анализ}
В силу симметрии маргинальных распределений $\tau_1$ и $\tau_2$, матрица ковариации является симметричной, то есть $c_{21}(t_1, t_2) = c_{12}(t_1, t_2)$. Итоговая ковариационная матрица процесса $D(t)$ имеет вид:
\begin{equation}
    C(t_1, t_2) = 
    \begin{pmatrix}
        c_{11}(t_1, t_2) & c_{12}(t_1, t_2) \\
        c_{12}(t_1, t_2) & c_{11}(t_1, t_2)
    \end{pmatrix},
\end{equation}
где элементы заданы формулами, выведенными выше.

\begin{remark}
Важно проанализировать поведение матрицы ковариации при $\alpha=0$ (случай независимых времен обслуживания внутри группы). При подстановке $\alpha=0$ в формулу (\ref{eq:cov_matrix_offdiag}) внедиагональные элементы не обращаются в ноль:
\begin{equation*}
    c_{12}(t_1, t_2) \Big|_{\alpha=0} = \lambda \left[ t_1 - 1 + e^{-t_1} + e^{-t_2} - \frac{1}{2}e^{t_1-t_2} - \frac{1}{2}e^{-t_1-t_2} \right] > 0.
\end{equation*}
Это аналитически доказывает фундаментальное свойство исследуемой системы: даже при независимых временах обслуживания ($\alpha=0$), сам факт одновременного поступления требований в систему (в составе одной группы) порождает строгую положительную ковариацию между компонентами выходящего процесса $D(t)$. Наличие параметра $\alpha \neq 0$ лишь дополнительно модифицирует эту зависимость.
\end{remark}

Чтобы понять физический и вероятностный смысл полученных формул, зафиксируем один и тот же момент времени ($t_1 = t_2 = t$) и исследуем элементы ковариационной матрицы выходящего процесса. На Рис. \ref{fig:cov_matrix} представлена динамика роста дисперсии и ковариации при интенсивности входящего потока $\lambda = 1$(соответствующий программный код приведен в Приложении \ref{app:python_code}).

\begin{figure}[h]
    \centering
    \includegraphics[width=\textwidth]{Images/covariance_matrix.pdf}
    \caption{Динамика диагональных (слева) и внедиагональных (справа) элементов ковариационной матрицы выходящего потока $D(t)$. Пунктирными линиями показаны линейные асимптоты процессов при $t \to \infty$.}
    \label{fig:cov_matrix}
\end{figure}

Анализ графиков и полученных асимптотик позволяет сделать следующие выводы о поведении системы:

\begin{enumerate}
    \item \textbf{Смысл ковариации в момент $t$.} Процесс $D_i(t)$ представляет собой счетчик требований типа $i$, покинувших систему к моменту $t$. Строгая положительность ковариации $c_{12}(t,t) > 0$ (см. правый график, $\alpha=0$) означает вероятностную синхронизацию объемов выполненной работы. Если к моменту $t$ система успела обслужить  большое число требований первого типа, она с высокой вероятностью выдаст пропорционально большое число требований второго типа, так как их поступление обусловлено общими группами пуассоновского процесса $A(t)$.
    
    \item \textbf{Асимптотика и пуассоновский предел.} Из формул (5) и (7) следует, что при $t \to \infty$ дисперсия и ковариация имеют линейные асимптоты:
    $$ Var(D_i(t)) = \lambda t - \lambda + o(1), $$
    $$ Cov(D_1(t), D_2(t)) = \lambda t - 1.5\lambda + \lambda \frac{\alpha}{12} + o(1). $$
    Линейный рост характеристик со скоростью $\lambda t$ является классическим индикатором того, что процесс уходов $D(t)$ «забывает» начальное нулевое состояние системы и выходит на режим со стационарными пуассоновскими приращениями, что согласуется с теоремой Бёрке \cite{burke} для стационарного распределения выходящего потока.
    
    \item \textbf{Асимптотическая полная корреляция потоков.} Самым нетривиальным результатом является поведение коэффициента корреляции $\rho_{12}(t)$ при $t \to \infty$. Используя полученные асимптотики, найдем предел коэффициента корреляции:
    $$ \lim_{t \to \infty} \rho_{12}(t) = \lim_{t \to \infty} \frac{Cov(D_1(t), D_2(t))}{\sqrt{Var(D_1(t)) \cdot Var(D_2(t))}} = \lim_{t \to \infty} \frac{\lambda t - 1.5\lambda + \lambda \frac{\alpha}{12}}{\sqrt{(\lambda t - \lambda)(\lambda t - \lambda)}} = 1. $$
    Данный предел строго доказывает своеобразный «закон сохранения» корреляции для рассматриваемой модели. Параметр внутренней зависимости времен обслуживания $\alpha \in [-1, 1]$ выступает лишь как локальная поправка (аддитивная константа), влияющая на дисперсию процессов на малых интервалах времени. Однако при $t \to \infty$ потоки обслуженных требований разных типов оказываются \textit{полностью скоррелированными} ($\rho = 1$) исключительно за счет факта их группового поступления.
\end{enumerate}

\section*{Заключение}
\addcontentsline{toc}{section}{Заключение}

В данной работе было проведено исследование выходящего потока $D(t)$ в системе массового обслуживания $M^k/G/\infty$ с групповым поступлением разнородных и зависимых требований. Использование вероятностного подхода, основанного на свойстве расщепления пуассоновского процесса, позволило преодолеть технические трудности классического метода производящих функций и получить интуитивно понятные, строгие результаты.

В ходе выполнения работы были получены следующие основные результаты:

\begin{enumerate}
    \item \textbf{Представление выходящего потока.} Доказана теорема, позволяющая выразить многомерный случайный процесс уходов $D(t)$ через линейную комбинацию независимых пуассоновских случайных величин $\delta_y(t)$. Это позволило свести сложную задачу анализа зависимых многомерных потоков к изучению независимых одномерных компонент.
    
    \item \textbf{Анализ ковариационных структур.} Разработан метод и получены общие интегральные представления для вычисления ковариационных матриц как между процессами числа требований в системе и выходящим потоком $\mathrm{Cov} (N(t), D(t))$, так и для автоковариационной функции самого выходящего процесса \protect \newline $\mathrm{Cov} (D(t_1), D(t_2))$.
    
    \item \textbf{Явные аналитические решения.} Для двумерного случая с использованием распределения Фарли-Гумбела-Моргенштерна (FGM) с экспоненциальными маргинальными распределениями получены точные аналитические формулы для элементов ковариационной матрицы выходящего потока. Продемонстрировано, как параметр внутренней зависимости времен обслуживания $\alpha$ влияет на дисперсию и ковариацию процессов на конечных интервалах времени.
    
    \item \textbf{Асимптотический анализ и «закон сохранения» корреляции.} Доказано, что при $t \to \infty$ дисперсия и ковариация процессов имеют линейные асимптоты, что согласуется с классической теоремой Бёрке о пуассоновской природе стационарного выходящего потока. Важнейшим вероятностным итогом работы является аналитическое доказательство того, что коэффициент корреляции потоков обслуженных требований разных типов стремится к единице:
    $$ \lim_{t \to \infty} \rho_{12}(t) = 1. $$
\end{enumerate}

Таким образом, в работе строго обосновано, что сам факт одновременного (группового) поступления заявок формирует сильную положительную зависимость в выходящем потоке. Внутренняя корреляция длительностей обслуживания внутри группы выступает лишь как локальное возмущение на малых масштабах времени, которое полностью исчезает на бесконечности, приводя потоки к полной асимптотической синхронизации. 

\newpage

\begin{thebibliography}{50}

\bibitem{choi1992}
Choi, B. D., Park, K. K. (1992). The $M^k/M/\infty$ queue with heterogeneous customers in a batch. \textit{Journal of applied probability}, 29(2), 477-481.

\bibitem{falin1994}
Falin, G. (1994). The $M^k/G/\infty$ batch arrival queue by heterogeneous dependent demands. \textit{Journal of applied probability}, 31(3), 841-846.

\bibitem{daw2022}
Daw, A., Fralix, B., Pender, J. (2020). Non-stationary queues with batch arrivals. arXiv preprint arXiv:2008.00625.

% \bibitem{nakamura2021}
% Nakamura, A., Phung-Duc, T. (2021, November). Stationary analysis of infinite server queue with batch service. In European Workshop on Performance Engineering (pp. 411-424). \textit{Cham: Springer International Publishing}.

% \bibitem{nakamura2024}
% Nakamura, A., Phung-Duc, T. (2024). Exact and asymptotic analysis of infinite server batch service queues with random batch sizes. \textit{Queueing Systems}, 106(1), 129-158.

\bibitem{gan2018}
Gan, Y., Delimitrou, C. (2018). The architectural implications of cloud microservices. \textit{IEEE Computer Architecture Letters}, 17(2), 155-158.

\bibitem{pang2012}
Pang, G.,  Whitt, W. (2012). Infinite-server queues with batch arrivals and dependent service times. \textit{Probability in the Engineering and Informational Sciences}, 26(2), 197-220.

\bibitem{Gumbel}
Gumbel, E. J. (1960). Bivariate exponential distributions. \textit{Journal of the American Statistical Association}, 55(292), 698-707.

% \bibitem{gubner2006}
% Gubner, J. A. (2006). Probability and random processes for electrical engineers. \textit{Cambridge University Press}.

\bibitem{burke}
Burke, P. J. (1956). The output of a queuing system. \textit{Operations research}, 4(6), 699-704.

\bibitem{bulinsky2005}
Булинский, А. В., Ширяев, А. Н. (2005). Теория случайных процессов. \textit{ФИЗМАТЛИТ}.

\bibitem{falinrisk}
Фалин Г.И. (1994). Математический анализ рисков в страховании. \textit{Российский юридический издательский дом}.

\bibitem{abolnikov1968}
Абольников, Л. М. (1968). Нестационарная задача массового обслуживания для систем с бесконечным числом каналов при групповом поступлении требовании. \textit{Проблемы передачи информации}, 4(3), 99-102.
\end{thebibliography}


\newpage
\appendix
\renewcommand{\thesection}{\Asbuk{section}}

\section*{Приложения}
\addcontentsline{toc}{section}{Приложения}

\titleformat{\section}{\normalfont\Large\bfseries}{Приложение \thesection.}{0.5em}{}

\section{Детальный разбор примера и визуализация процесса $\delta_y(t)$}
\label{app:example}

Чтобы наглядно продемонстрировать поведение случайного процесса $\delta_y(t)$, рассмотрим конкретный сценарий и пошагово рассчитаем его значения.

\subsubsection*{1. Сценарий}

Пусть в системе есть требования двух типов, т.е. $S = \{1, 2\}$. Мы хотим отслеживать количество процессов, находящихся в состоянии $y = \{1\}$, что означает: <<первое требование обслужено, а второе -- еще нет>>.

Рассмотрим поведение системы с тремя группами($n=1, 2, 3$), которые поступают в разное время и имеют разное время обслуживания требований, причем пока они все не обслужены, новых групп не поступает:
\begin{itemize}
    \item \textbf{Группа 1 ($n=1$):}
    \begin{itemize}
        \item Время поступления: $\alpha_1 = 1$.
        \item Время обслуживания: $\tau_1^{(1)} = 3$ (завершение в $t=1+3=4$), $\tau_2^{(1)} = 5$ (завершение в $t=1+5=6$).
    \end{itemize}
    \item \textbf{Группа 2 ($n=2$):}
    \begin{itemize}
        \item Время поступления: $\alpha_2 = 2$.
        \item Время обслуживания: $\tau_1^{(2)} = 2$ (завершение в $t=2+2=4$), $\tau_2^{(2)} = 6$ (завершение в $t=2+6=8$).
    \end{itemize}
    \item \textbf{Группа 3 ($n=3$):}
    \begin{itemize}
        \item Время поступления: $\alpha_3 = 5$.
        \item Время обслуживания: $\tau_1^{(3)} = 2$ (завершение в $t=5+2=7$), $\tau_2^{(3)} = 4$ (завершение в $t=5+4=9$).
    \end{itemize}
\end{itemize}

\subsubsection*{2. Таблица состояний процессов}

Проследим, как меняется множество обслуженных требований $Y_n(t)$ для каждой группы с течением времени. Целевое состояние $y=\{1\}$ выделено жирным шрифтом.

\begin{table}[H]
\centering
\caption{$Y_n(t)$ для каждой группы}
\renewcommand{\arraystretch}{1.2}
\begin{tabular}{|c|c|c|c|}
\hline
\textbf{Интервал времени ($t$)} & \textbf{$Y_1(t)$} & \textbf{$Y_2(t)$} & \textbf{$Y_3(t)$} \\ \hline
$[0, 4)$ & $\emptyset$ & $\emptyset$ & $\emptyset$ \\ \hline
$[4, 6)$ & $\mathbf{\{1\}}$ & $\mathbf{\{1\}}$ & $\emptyset$ \\ \hline
$[6, 7)$ & $\{1, 2\}$ & $\mathbf{\{1\}}$ & $\emptyset$ \\ \hline
$[7, 8)$ & $\{1, 2\}$ & $\mathbf{\{1\}}$ & $\mathbf{\{1\}}$ \\ \hline
$[8, 9)$ & $\{1, 2\}$ & $\{1, 2\}$ & $\mathbf{\{1\}}$ \\ \hline
$[9, +\infty)$ & $\{1, 2\}$ & $\{1, 2\}$ & $\{1, 2\}$ \\ \hline
\end{tabular}
\end{table}

\subsubsection*{3. Расчет значений $\delta_y(t)$}

Используя таблицу, посчитаем значение $\delta_{\{1\}}(t) = \sum_{n=1}^{3} \mathbb{I}(Y_n(t)=\{1\})$ для каждого интервала.
\begin{itemize}
    \item $0 \le t < 4$: Все группы в состоянии $\emptyset$. Значения индикаторов: $0, 0, 0$. \\
    Следовательно, $\delta_{\{1\}}(t) = 0 + 0 + 0 = 0$.

    \item $4 \le t < 6$: Группы 1 и 2 находятся в состоянии $\{1\}$, группа 3 -- в $\emptyset$. Значения индикаторов: $1, 1, 0$. \\
    Следовательно, $\delta_{\{1\}}(t) = 1 + 1 + 0 = 2$.

    \item $6 \le t < 7$: Группа 1 перешла в $\{1, 2\}$, группа 2 осталась в $\{1\}$. Значения индикаторов: $0, 1, 0$. \\
    Следовательно, $\delta_{\{1\}}(t) = 0 + 1 + 0 = 1$.

    \item $7 \le t < 8$: Группа 2 все еще в $\{1\}$, Группа 3 вошла в это состояние. Значения индикаторов: $0, 1, 1$. \\
    Следовательно, $\delta_{\{1\}}(t) = 0 + 1 + 1 = 2$.

    \item $8 \le t < 9$: Группа 2 покинула $\{1\}$, Группа 3 осталась. Значения индикаторов: $0, 0, 1$. \\
    Следовательно, $\delta_{\{1\}}(t) = 0 + 0 + 1 = 1$.

    \item $t \ge 9$: Все группы покинули состояние $\{1\}$. Значения индикаторов: $0, 0, 0$. \\
    Следовательно, $\delta_{\{1\}}(t) = 0 + 0 + 0 = 0$.
\end{itemize}

Этот пошаговый разбор показывает, как формируется ступенчатый график, представленный ниже.

\begin{figure}[H]
\centering
\begin{tikzpicture}
\pgfplotsset{compat=1.18} 

\begin{axis}[
    xlabel={$t$},
    xmin=0, xmax=12,
    ymin=0, ymax=3,
    axis x line=bottom,
    axis y line=left,
    ytick={0,1,2,3},
    grid=major,
    grid style={dashed, gray!30},
    width=13cm,
    height=7cm,
]

\draw[blue, ultra thick] (axis cs:0,0) -- (axis cs:4,0);
\draw[blue, thick, dashed] (axis cs:4,0) -- (axis cs:4,2);

\draw[blue, ultra thick] (axis cs:4,2) -- (axis cs:6,2);
\draw[blue, thick, dashed] (axis cs:6,2) -- (axis cs:6,1);

\draw[blue, ultra thick] (axis cs:6,1) -- (axis cs:7,1);
\draw[blue, thick, dashed] (axis cs:7,1) -- (axis cs:7,2);

\draw[blue, ultra thick] (axis cs:7,2) -- (axis cs:8,2);
\draw[blue, thick, dashed] (axis cs:8,2) -- (axis cs:8,1);

\draw[blue, ultra thick] (axis cs:8,1) -- (axis cs:9,1);
\draw[blue, thick, dashed] (axis cs:9,1) -- (axis cs:9,0);

\draw[blue, ultra thick] (axis cs:9,0) -- (axis cs:12,0);

\addplot+[blue, only marks, mark=*, mark options={fill=blue}] coordinates {
    (0, 0)
    (4, 2)
    (6, 1)
    (7, 2)
    (8, 1)
    (9, 0)
    (12, 0)
};
\addlegendentry{$\delta_{\{1\}}(t)$}

\end{axis}
\end{tikzpicture}
\caption{График $\delta_y(t)$ для $y = \{1\}$}
\label{fig:delta_plot_custom_style}
\end{figure}




\newpage
\section{Исходный код на языке Python для генерации графиков}
\label{app:python_code}

В данном приложении представлены программы на языке Python, использованные для численного моделирования и построения графиков (Рис. \ref{fig:fgm_contours}, \ref{fig:g1_plot} и \ref{fig:cov_matrix}). Графики построены с использованием библиотек \texttt{matplotlib} и \texttt{numpy}.

\subsection*{Код для построения линий уровня (Рис. \ref{fig:fgm_contours})}
\begin{lstlisting}
import numpy as np
import matplotlib.pyplot as plt

plt.rcParams.update({
    "font.family": "serif",
    "mathtext.fontset": "cm",
    "axes.labelsize": 14,
    "xtick.labelsize": 12,
    "ytick.labelsize": 12,
    "axes.titlesize": 14
})

def fgm_pdf(t1, t2, alpha):
    return np.exp(-t1 - t2) * (1 + alpha * (2 * np.exp(-t1) - 1) * (2 * np.exp(-t2) - 1))

t1 = np.linspace(0.01, 4, 150)
t2 = np.linspace(0.01, 4, 150)
T1, T2 = np.meshgrid(t1, t2)

alphas = [-1, 0, 1]
titles = [
    r"$\alpha = -1$ (Отрицательная)",
    r"$\alpha = 0$ (Независимость)",
    r"$\alpha = 1$ (Положительная)"
]

fig, axes = plt.subplots(1, 3, figsize=(15, 5))

for ax, alpha, title in zip(axes, alphas, titles):
    Z = fgm_pdf(T1, T2, alpha)
    
    contour_filled = ax.contourf(T1, T2, Z, levels=25, cmap='Blues')
    
    ax.contour(T1, T2, Z, levels=25, colors='black', linewidths=0.3, alpha=0.5)
    
    ax.set_title(title, pad=15)
    ax.set_xlabel(r"Время обслуживания $t_1$")
    ax.set_ylabel(r"Время обслуживания $t_2$")
    ax.grid(True, linestyle='--', alpha=0.4)
    ax.set_aspect('equal') 

plt.tight_layout()
plt.savefig("fgm_density_contours.pdf", bbox_inches='tight', dpi=300)
plt.show()
\end{lstlisting}

\subsection*{Код для динамики функции $g_{\{1\}}(t)$ (Рис. \ref{fig:g1_plot})}
\begin{lstlisting}
import numpy as np
import matplotlib.pyplot as plt

plt.rcParams.update({
    "font.family": "serif",
    "mathtext.fontset": "cm",
    "font.size": 12
})

t = np.linspace(0, 10, 500)

def g1(t, alpha):
    term1 = 1 - np.exp(-t)
    term2 = ((alpha + 1) / 2) * (1 - np.exp(-2*t))
    term3 = (2 * alpha / 3) * (1 - np.exp(-3*t))
    term4 = (alpha / 4) * (1 - np.exp(-4*t))
    return term1 - term2 + term3 - term4

plt.figure(figsize=(9, 5.5))

plt.plot(t, g1(t, -1), color='red', linestyle='-', linewidth=2, 
         label=r'$\alpha = -1$ (отрицательная зависимость)')
plt.plot(t, g1(t, 0), color='black', linestyle='-', linewidth=2, 
         label=r'$\alpha = 0$ (независимость)')
plt.plot(t, g1(t, 1), color='blue', linestyle='-', linewidth=2, 
         label=r'$\alpha = 1$ (положительная зависимость)')

plt.axhline(y=7/12, color='red', linestyle='--', alpha=0.6)   
plt.axhline(y=1/2, color='black', linestyle='--', alpha=0.6)  
plt.axhline(y=5/12, color='blue', linestyle='--', alpha=0.6)  

plt.yticks([0, 0.2, 5/12, 1/2, 7/12, 0.7], 
           ['0', '0.2', r'$5/12$', r'$1/2$', r'$7/12$', '0.7'])

plt.title(r'Динамика функции $g_{\{1\}}(t)$ и выход на стационарный режим', fontsize=14, pad=15)
plt.xlabel(r'Время $t$', fontsize=12)
plt.ylabel(r'Значение $g_{\{1\}}(t)$', fontsize=12)

plt.grid(True, linestyle=':', alpha=0.7)
plt.legend(loc='lower right', fontsize=11, framealpha=0.9)

plt.xlim(0, 7)
plt.ylim(0, 0.65)

plt.savefig('g1_plot.pdf', bbox_inches='tight')
print("График успешно сохранен как g1_plot.pdf")
\end{lstlisting}

\subsection*{Код для элементов ковариационной матрицы (Рис. \ref{fig:cov_matrix})}
\begin{lstlisting}
import numpy as np
import matplotlib.pyplot as plt

plt.rcParams.update({
    "font.family": "serif",
    "mathtext.fontset": "cm",
    "axes.labelsize": 14,
    "xtick.labelsize": 12,
    "ytick.labelsize": 12,
    "axes.titlesize": 14,
    "legend.fontsize": 12
})

t = np.linspace(0, 8, 400)
lmbda = 1  

c11 = lmbda * (t - 1 + np.exp(-t))
c11_asymp = lmbda * (t - 1)  

def c12(t, alpha):
    base = t - 1.5 + 2 * np.exp(-t) - 0.5 * np.exp(-2*t)
    correction = alpha * (1/12 + (1/6) * np.exp(-2*t) - (1/4) * np.exp(-4*t))
    return lmbda * (base + correction)

fig, axes = plt.subplots(1, 2, figsize=(14, 6))

axes[0].plot(t, c11, 'b-', linewidth=2, label=r'Дисперсия $c_{11}(t,t)$')
axes[0].plot(t, c11_asymp, 'k--', linewidth=1.5, alpha=0.7, label=r'Асимптота $\lambda(t-1)$')
axes[0].set_title(r'Диагональный элемент матрицы (Дисперсия $c_{11}(t,t)$)')
axes[0].set_xlabel(r'Время $t$')
axes[0].set_ylabel(r'Значение')
axes[0].grid(True, linestyle='--', alpha=0.6)
axes[0].legend(loc='upper left')
axes[0].set_xlim([0, 5])
axes[0].set_ylim([0, 5])


axes[1].plot(t, c12(t, 1), 'r-', linewidth=2, label=r'$\alpha = 1$ (Положит. завис.)')
axes[1].plot(t, c12(t, 0), 'g-', linewidth=2, label=r'$\alpha = 0$ (Независимость)')
axes[1].plot(t, c12(t, -1), 'b-', linewidth=2, label=r'$\alpha = -1$ (Отрицат. завис.)')

axes[1].plot(t, t - 1.5 + 1/12, 'r--', linewidth=1, alpha=0.5)
axes[1].plot(t, t - 1.5, 'g--', linewidth=1, alpha=0.5)
axes[1].plot(t, t - 1.5 - 1/12, 'b--', linewidth=1, alpha=0.5)

axes[1].set_title(r'Внедиагональный элемент (Ковариация $c_{12}(t,t)$)')
axes[1].set_xlabel(r'Время $t$')
axes[1].set_ylabel(r'Значение ковариации')
axes[1].grid(True, linestyle='--', alpha=0.6)
axes[1].legend(loc='upper left')
axes[1].set_xlim([0, 5])
axes[1].set_ylim([0, 5])

plt.tight_layout()
plt.savefig("covariance_matrix.pdf", bbox_inches='tight', dpi=300)
plt.show()
\end{lstlisting}

\end{document}
